\documentclass[11pt,a4paper]{article}
\usepackage{notes}
\usepackage{fancyhdr}

\pagestyle{fancy}
\fancyhf{}
\lhead{Geophysics 3A: Potential Fields \& Geothermics}
\rhead{Week 7 Fourier Transforms Demo}
\cfoot{\thepage}

\title{Week 7 --- Demo: Fourier Decomposition: Amplitude, Phase, and Sampling}
\author{Geophysics 3A: Potential Fields \& Geothermics}
\date{\today}

\begin{document}
\maketitle

\section{Purpose of the Demonstration}

This demonstration is designed to build intuition for what a Fourier transform does and does not do.
In particular, it addresses:

\begin{itemize}
  \item How a signal may be constructed from known frequencies, amplitudes, and phases
  \item How a discrete Fourier transform (DFT) estimates those parameters from sampled data
  \item Why amplitude recovery is affected by sampling and windowing
  \item Why phase recovery depends on reference conventions
  \item How these ideas carry directly into spatial spectral methods in potential field geophysics
\end{itemize}

Although the demonstration uses time-domain signals, the mathematics applies identically to spatial data.

\section{Signal Construction}

We construct a signal as a sum of sinusoidal components:
\begin{equation}
x(t) = \sum_{i=1}^{N_f} A_i \sin\!\left(2\pi f_i t + \phi_i\right),
\end{equation}
where:
\begin{itemize}
  \item $f_i$ are known frequencies,
  \item $A_i$ are amplitudes,
  \item $\phi_i$ are phase offsets.
\end{itemize}

In the demonstration, amplitudes follow a power-law decay,
\begin{equation}
A_i = f_i^{-n},
\end{equation}
to mimic the tendency for geophysical spectra to be dominated by long wavelengths.

\subsection{Linear Phase and Time Shifts}

A linear phase is imposed by choosing
\begin{equation}
\phi_i = -2\pi f_i t_0,
\end{equation}
which gives
\begin{equation}
x(t) = \sum_i A_i \sin\!\left[2\pi f_i (t - t_0)\right].
\end{equation}

Thus, a linear phase spectrum corresponds to a constant time shift $t_0$.
This is a fundamental Fourier property:
\begin{equation}
x(t - t_0) \;\Longleftrightarrow\; X(f)\,e^{-i 2\pi f t_0}.
\end{equation}

The key information carried by phase in this context is therefore the \emph{slope} of phase as a function of frequency, not its absolute value.

\section{Discrete Fourier Transform}

A discrete Fourier transform of $N$ samples $x_n = x(n\Delta t)$ is defined as
\begin{equation}
X_k = \sum_{n=0}^{N-1} x_n \, e^{-i 2\pi kn/N},
\end{equation}
where $k$ indexes discrete frequency bins
\begin{equation}
f_k = \frac{k}{N\,\Delta t}.
\end{equation}

The FFT returns complex coefficients $X_k$; from these we extract:
\begin{align}
\text{Amplitude: } & |X_k|, \\
\text{Phase: } & \arg(X_k).
\end{align}

\subsection{Amplitude Scaling}

The DFT does not return amplitudes in the same units used to construct the signal.
For a pure sinusoid of amplitude $A$ sampled $N$ times, the magnitude of its FFT coefficient scales as $\sim N/2$.

To recover physical amplitudes from a one-sided spectrum, we apply the deterministic scaling
\begin{equation}
A_k^{\text{(rec)}} = \frac{2}{N\,G_w} |X_k|,
\end{equation}
where $G_w$ is the coherent gain of the window function.
For a Hann window,
\begin{equation}
G_w = \frac{1}{N} \sum_{n=0}^{N-1} w_n \approx 0.5.
\end{equation}

This correction removes known numerical scaling but does not make the recovery exact; remaining discrepancies arise from finite record length, bin mismatch, and spectral leakage.

\section{Phase Conventions and Interpretation}

Phase recovery requires particular care.

\subsection{Sine versus Cosine Basis}

The signal is constructed using sine functions, while the FFT basis corresponds to complex exponentials (or equivalently, cosines and sines).
Since
\begin{equation}
\sin(\omega t + \phi) = \cos\!\left(\omega t + \phi - \frac{\pi}{2}\right),
\end{equation}
the recovered FFT phase includes a constant offset of $-\pi/2$ relative to the input phase.
This offset is frequency-independent and carries no physical meaning.

\subsection{Absolute Phase versus Phase Differences}

Phase is inherently defined modulo $2\pi$, and an overall constant phase depends on the choice of time origin.
In this demonstration:
\begin{itemize}
  \item Absolute phase values depend on arbitrary reference choices.
  \item Phase \emph{differences} and \emph{slopes} carry physical information.
\end{itemize}

This distinction mirrors geophysical practice:
\begin{itemize}
  \item In magnetotellurics, phase is referenced to the inducing field and is physically interpretable.
  \item In potential fields and this demo, constant phase offsets correspond to arbitrary spatial or temporal origins.
\end{itemize}

\section{Connection to Spatial Geophysics}

The mathematics is identical if time $t$ is replaced by spatial position $x$:
\begin{table}
    \centering
    \begin{tabular}{ll}
        \toprule
        \textbf{Time-domain signal} &
        \textbf{Spatial potential field} \\
        \midrule
        $t$ & $x$ \\
        Frequency $f$ & Wavenumber $k$ \\
        Time shift & Lateral offset \\
        Sampling rate & Station spacing \\
        Linear phase & Spatial translation \\
        \bottomrule
    \end{tabular}
\end{table}

These ideas underpin upward and downward continuation, reduction to the pole, and spectral depth estimation.

\section{Suggested In-Class Demonstration Outline}

The following sequence uses the interactive notebook widgets.

\subsection*{Step 1: Frequency Recovery}
\begin{itemize}
  \item Phase type: linear
  \item Phase shift: $t_0 \approx 0.2$ s
  \item Amplitude exponent: $n = 1$
  \item Padding: $\times 4$
  \item Window: on
\end{itemize}
Observe that FFT peaks occur near the input frequencies.

\subsection*{Step 2: Amplitude Decay}
Increase the amplitude exponent $n$ from 0 to 2 and observe how spectral shape changes while phase trends persist.

\subsection*{Step 3: Linear Phase and Time Shifts}
Vary $t_0$ and observe:
\begin{itemize}
  \item A linear phase trend
  \item Increasing slope with increasing shift
  \item No change in amplitude spectrum
\end{itemize}

\subsection*{Step 4: Phase Scrambling}
Set phase type to ``scrambled.''
Observe that:
\begin{itemize}
  \item The time series becomes complex
  \item The amplitude spectrum is unchanged
  \item Phase loses coherent structure
\end{itemize}

\subsection*{Step 5: Padding and Windowing}
Toggle padding and windowing to show:
\begin{itemize}
  \item Improved phase stability with padding
  \item Reduced leakage with windowing
\end{itemize}

Emphasize that improvements are numerical, not physical.

\section{Key Takeaways}

\begin{itemize}
  \item Fourier transforms reorganize information; they do not interpret it.
  \item Amplitude and phase are recovered approximately from sampled data.
  \item Phase slopes carry physical meaning; absolute offsets often do not.
  \item The same concepts govern spectral methods in potential field geophysics.
\end{itemize}

\end{document}